\chapter{Espaces de Hilbert}\label{ch:b3:hilbert}

Un \omterm{def:b3:hilbert:inner}{espace de Hilbert} est un espace de Banach dont la norme
provient d'un \omterm{def:b3:hilbert:inner}{produit scalaire} --- et cette unique structure
supplémentaire restaure, en dimension infinie, presque toute la
géométrie euclidienne~: les \omterm{thm:b3:hilbert:decomposition}{projections orthogonales} existent,
toute forme linéaire \omterm{def:b3:topology:continuity}{continue} est un \omterm{def:b3:hilbert:inner}{produit scalaire} contre un
vecteur fixe (Riesz), et les \omterm{def:b3:topology:basis}{bases} orthonormées développent
chaque vecteur en une série convergente avec la comptabilité
pythagoricienne (\omterm{thm:b3:hilbert:parseval}{Parseval}). Le point culminant du chapitre
honore une dette du volume de deuxième année~: le système
trigonométrique est une \omterm{def:b3:hilbert:onb}{base hilbertienne} de $L^2$, si bien que
l'identité de \omterm{thm:b3:hilbert:parseval}{Parseval} vaut pour \emph{toute} fonction de carré
\omterm{def:b3:lebesgue:l1}{intégrable} --- l'énoncé que la deuxième année ne pouvait
démontrer que pour les fonctions $\mathcal C^1$ par morceaux.
Nous terminons par Lax--Milgram, le lemme de travail de
l'approche variationnelle des équations différentielles.

Tout au long du chapitre, $H$ est un espace vectoriel sur
$K = \R$ ou $\C$.

\section{Produits scalaires~; le théorème de projection}

\begin{definition}\label{def:b3:hilbert:inner}
Un \emph{produit scalaire}\index{produit scalaire} est une
application $\langle \cdot,\cdot\rangle \colon H\times H \to
K$, linéaire en la seconde variable, avec $\langle y, x\rangle
= \overline{\langle x, y\rangle}$ et $\langle x, x\rangle > 0$
pour $x \neq 0$. Elle induit la norme $\norm x = \langle x,
x\rangle^{1/2}$, l'\emph{inégalité de Cauchy--Schwarz}
$\abs{\langle x, y\rangle} \leq \norm x\norm y$ (la preuve de
deuxième année --- le discriminant --- est inchangée), et la
\emph{identité du parallélogramme}
\[
\norm{x + y}^2 + \norm{x - y}^2 = 2\norm x^2 + 2\norm y^2 .
\]
Un \emph{espace de Hilbert}\index{espace de Hilbert} est un
espace préhilbertien \omterm{def:b3:complete:complete}{complet} pour cette norme. Exemples~:
$\ell^2$ (\cref{pb:b3:banach:1}) et, l'exemple fondamental,
$L^2(\mu)$ avec $\langle f, g\rangle = \int\bar fg\,\dd\mu$
--- \omterm{def:b3:complete:complete}{complet} par Riesz--Fischer (\cref{thm:b3:lp:complete})~;
le produit scalaire est fini par Cauchy--Schwarz ($=$ Hölder
en $p = q = 2$).
\end{definition}

\begin{theorem}[Projection sur un convexe fermé]
\label{thm:b3:hilbert:projection}
Soit $C \neq \varnothing$ un sous-ensemble \emph{convexe}
fermé de l'\omterm{def:b3:hilbert:inner}{espace de Hilbert} $H$ et $x \in H$. Il existe un
unique $p_C(x) \in C$ tel que\index{théorème de projection}
\[
\norm{x - p_C(x)} = d(x, C),
\]
caractérisé par~: $\operatorname{Re}\langle x - p_C(x),\ c -
p_C(x)\rangle \leq 0$ pour tout $c \in C$. L'application
$p_C$ est $1$-lipschitzienne.
\end{theorem}

\begin{proof}
Soit $d = d(x, C)$ et $(c_n) \subseteq C$ avec $\norm{x -
c_n} \to d$. Parallélogramme sur $x - c_n$ et $x - c_m$~:
\[
\norm{c_n - c_m}^2 = 2\norm{x - c_n}^2 + 2\norm{x - c_m}^2 -
4\,\bigl\|x - \tfrac{c_n + c_m}2\bigr\|^2
\leq 2\norm{x{-}c_n}^2 + 2\norm{x{-}c_m}^2 - 4d^2
\]
(la convexité place le milieu dans $C$)~: le membre de droite
tend vers $0$, donc $(c_n)$ est de Cauchy, et sa limite $p
\in C$ (fermé) réalise $d$. Unicité~: deux minimiseurs
donnent, par la même identité, $\norm{p - p'}^2 \leq 2d^2 +
2d^2 - 4d^2 = 0$.

Caractérisation~: pour $c \in C$, $t \in \intoc01$, le
vecteur $p + t(c - p) \in C$, d'où
\[
d^2 \leq \norm{x - p - t(c-p)}^2
= d^2 - 2t\operatorname{Re}\langle x - p, c - p\rangle +
t^2\norm{c-p}^2 ;
\]
on divise par $t \to 0^+$~: $\operatorname{Re}\langle x - p,
c - p\rangle \leq 0$. Réciproquement cette inégalité donne
$\norm{x - c}^2 = \norm{x - p}^2 - 2\operatorname{Re}\langle
x - p, c - p\rangle + \norm{p - c}^2 \geq \norm{x-p}^2$.
Lipschitz~: pour $x, y$ de projections $p, q$, on additionne
les deux inégalités variationnelles (avec $c = q$, resp.\
$c = p$)~: $\operatorname{Re}\langle x - y - (p - q), p -
q\rangle \geq 0$, d'où $\norm{p - q}^2 \leq
\operatorname{Re}\langle x - y, p - q\rangle \leq \norm{x -
y}\norm{p - q}$.
\end{proof}

\begin{theorem}[Décomposition orthogonale]
\label{thm:b3:hilbert:decomposition}
Soit $F$ un \emph{sous-espace fermé} de $H$. Alors $p_F$ est
linéaire, $x - p_F(x) \perp F$ pour tout $x$, et\index{orthogonal
complement}
\[
H = F \oplus F^\perp,
\qquad F^\perp = \{y : \langle y, f\rangle = 0\ \forall f\in
F\},
\qquad (F^\perp)^\perp = F .
\]
Pour un sous-espace quelconque, $(F^\perp)^\perp = \bar F$~; en
particulier $F$ est \omterm{def:b3:topology:interior}{dense} ssi $F^\perp = \{0\}$.
\end{theorem}

\begin{proof}
Pour un sous-espace, la caractérisation variationnelle avec $c
= p_F(x) \pm f$ ($f \in F$, les deux signes, et $\iu f$ dans le
cas complexe) force $\langle x - p_F(x), f\rangle = 0$~: le
résidu est orthogonal à $F$. Décomposition $x = p_F(x) + (x
- p_F(x))$ avec $F \cap F^\perp = \{0\}$ ($\langle y, y
\rangle = 0$)~; la linéarité de $p_F$ suit de l'unicité de
telles décompositions (les deux côtés y sont linéaires).
$(F^\perp)^\perp \supseteq F$ toujours~; réciproquement si $x
\perp F^\perp$, on écrit $x = f + g$~: $g = x - f \in
F^\perp$ et $\langle g, g\rangle = \langle x, g\rangle -
\langle f, g\rangle = 0$~: $x = f \in F$. Pour un sous-espace
général $F$~: $F^\perp = \bar F^{\,\perp}$ (continuité du
\omterm{def:b3:hilbert:inner}{produit scalaire}), donc $(F^\perp)^\perp = \bar F$ par le cas
fermé~; \omterm{ex:b3:lebesgue:gamma}{densité} ssi $\bar F = H$ ssi $F^\perp = 0$.
\end{proof}

\begin{theorem}[Représentation de Riesz]\label{thm:b3:hilbert:riesz}
Pour toute forme linéaire \omterm{def:b3:topology:continuity}{continue} $\varphi \in H'$ il
existe un unique $a \in H$ tel que\index{théorème de
représentation de Riesz}
\[
\varphi(x) = \langle a, x\rangle \quad (x \in H),
\qquad \norm\varphi_{H'} = \norm a .
\]
\end{theorem}

\begin{proof}
Si $\varphi = 0$~: $a = 0$. Sinon $F = \ker\varphi$ est un
sous-espace fermé propre~; on prend $u \in F^\perp$, $\norm u
= 1$ (\cref{thm:b3:hilbert:decomposition}~: $F^\perp \neq 0$
puisque $F \neq H$). Pour tout $x$, le vecteur $\varphi(x)u -
\varphi(u)x \in \ker\varphi$, donc $\perp u$~:
\[
0 = \langle u, \varphi(x)u - \varphi(u)x\rangle
= \varphi(x) - \varphi(u)\langle u, x\rangle :
\qquad \varphi(x) = \langle
\overline{\varphi(u)}\,u,\ x\rangle .
\]
Donc $a = \overline{\varphi(u)}u$ convient. Unicité~:
$\langle a - a', x\rangle = 0$ pour tout $x$, on teste $x = a
- a'$. Normes~: $\abs{\varphi(x)} \leq \norm a\norm x$
(Cauchy--Schwarz) avec égalité en $x = a$.
\end{proof}

\begin{example}[Une projection, calculée jusqu'au bout]
\label{ex:b3:hilbert:projectionexample}
Dans $H = L^2(\intcc01)$, quelle est la meilleure
approximation de $f(x) = x^2$ par une fonction affine~? Le
sous-espace $F = \operatorname{Vect}(1, x)$ est fermé
(dimension finie), et $p_F(f) = a + bx$ est caractérisé par
l'orthogonalité du résidu à $1$ et à $x$~:
\[
\int_0^1(x^2 - a - bx)\,\dd x = 0,
\qquad
\int_0^1x\,(x^2 - a - bx)\,\dd x = 0,
\]
c'est-à-dire $\frac13 = a + \frac b2$ et $\frac14 = \frac a2
+ \frac b3$~: $a = -\frac16$, $b = 1$. Donc $p_F(x^2) = x -
\frac16$, et l'erreur vaut
\[
d(f, F)^2 = \int_0^1\Bigl(x^2 - x + \frac16\Bigr)^2\dd x =
\frac1{180},
\qquad d(f, F) = \frac1{6\sqrt5} .
\]
Deux remarques à intérioriser. D'abord, le calcul n'est rien
d'autre qu'un système linéaire $2\times2$ --- les
\emph{équations \omterm{def:b3:groups:normal}{normales}}~; pour la \omterm{def:b3:topology:basis}{base} monomiale leur
matrice $\bigl(\frac1{i+j+1}\bigr)$ est la tristement célèbre
matrice de Hilbert, mal conditionnée, et orthonormaliser
d'abord (polynômes de Legendre, \cref{pb:b3:hilbert:1}) est
le remède. Ensuite, la meilleure approximation
\emph{uniforme} de $x^2$ par des affines est différente ($x
- \frac18$, par équi-oscillation)~: chaque norme a sa
géométrie, et seule la géométrie hilbertienne répond par un
système linéaire.
\end{example}

\section{Bases orthonormées}

\begin{definition}\label{def:b3:hilbert:onb}
Une famille $(e_i)_{i\in I}$ est \emph{orthonormée} si
$\langle e_i, e_j\rangle = \delta_{ij}$, et une \emph{base
hilbertienne}\index{base hilbertienne} (\omterm{def:b3:topology:basis}{base} orthonormée) si
de plus ses combinaisons linéaires finies sont \omterm{def:b3:topology:interior}{denses} dans
$H$ (la famille est \emph{totale}). On traite le cas dénombrable
$I = \N$, qui par Gram--Schmidt couvre tout $H$
\emph{\omterm{def:b3:galois:separable}{séparable}} (\cref{prop:b3:hilbert:gramschmidt}).
\end{definition}

\begin{theorem}[Bessel, Parseval]\label{thm:b3:hilbert:parseval}
Soit $(e_n)_{n\in\N}$ orthonormée dans $H$, et $c_n(x) =
\langle e_n, x\rangle$.
\begin{enumerate}
  \item (Bessel) $\sum_n\abs{c_n(x)}^2 \leq \norm x^2$, et
        la série $\sum_nc_n(x)e_n$ converge dans $H$, de
        somme $p_F(x)$, $F = \overline{\operatorname{Vect}}(e_n)$.
  \item Les assertions suivantes sont équivalentes~: (i)
        $(e_n)$ est une \omterm{def:b3:hilbert:onb}{base hilbertienne}~; (ii) $x =
        \sum_nc_n(x)e_n$ pour tout $x$~; (iii)
        \emph{Parseval}\index{identité de Parseval}~: $\norm
        x^2 = \sum_n\abs{c_n(x)}^2$ pour tout $x$~; (iv) le
        seul vecteur orthogonal à tous les $e_n$ est $0$.
  \item Si $(e_n)$ est une \omterm{def:b3:hilbert:onb}{base hilbertienne}, $x \mapsto
        (c_n(x))_n$ est un isomorphisme isométrique $H \to
        \ell^2$ (\emph{tout} \omterm{def:b3:hilbert:inner}{espace de Hilbert} \omterm{def:b3:galois:separable}{séparable} de
        dimension infinie «~est~» $\ell^2$), et $\langle x,
        y\rangle = \sum_n\overline{c_n(x)}c_n(y)$.
\end{enumerate}
\end{theorem}

\begin{proof}
(1) Pour $N$ fini~: $x - \sum_{n\leq N}c_ne_n \perp e_k$ ($k
\leq N$), donc Pythagore donne $\norm x^2 = \sum_{n\leq
N}\abs{c_n}^2 + \norm{x - \sum_{n\leq N}c_ne_n}^2$~: Bessel.
Les sommes partielles $S_N = \sum_{n\leq N}c_ne_n$ sont de
Cauchy~: $\norm{S_N - S_M}^2 = \sum_{M<n\leq N}\abs{c_n}^2$,
queue d'une série convergente~; la limite est dans $F$, et $x
- \lim S_N \perp$ chaque $e_k$ (continuité), donc $\perp F$~:
par unicité de la \omterm{thm:b3:hilbert:decomposition}{décomposition orthogonale}, $\lim S_N =
p_F(x)$.

(2) (i)$\Rightarrow$(ii)~: $F = H$, donc $p_F = \mathrm{id}$.
(ii)$\Rightarrow$(iii)~: Pythagore à la limite ($\norm{S_N}^2
= \sum_{n \leq N}\abs{c_n}^2 \to \norm x^2$).
(iii)$\Rightarrow$(iv)~: $x \perp$ tous les $e_n$ donne
$\norm x^2 = 0$. (iv)$\Rightarrow$(i)~: $F^\perp = \{0\}$
(l'orthogonalité à tous les $e_n$ est l'orthogonalité à $F$),
donc $F$ est \omterm{def:b3:topology:interior}{dense} par
le~\cref{thm:b3:hilbert:decomposition}~; mais $F$, une
fermeture, est déjà fermé~: $F = H$.

(3) L'application est linéaire, isométrique par (iii) (donc
injective), et surjective~: pour $(c_n) \in \ell^2$, la série
$\sum c_ne_n$ converge (Cauchy comme en (1)) vers un
antécédent. La formule du \omterm{def:b3:hilbert:inner}{produit scalaire} est la polarisation
de (iii), ou un calcul direct de limite.
\end{proof}

\begin{proposition}[Gram--Schmidt]\label{prop:b3:hilbert:gramschmidt}
Soit $(x_n)$ une suite linéairement indépendante. En posant
inductivement $\tilde e_n = x_n - \sum_{k<n}\langle e_k,
x_n\rangle e_k$ et $e_n = \tilde e_n/\norm{\tilde e_n}$ on
produit une famille orthonormée $(e_n)$ de mêmes enveloppes
finies~: $\operatorname{Vect}(e_1, \dots, e_n) =
\operatorname{Vect}(x_1, \dots, x_n)$. Par suite tout \omterm{def:b3:hilbert:inner}{espace
de Hilbert} \omterm{def:b3:galois:separable}{séparable} (muni d'une \omterm{def:b3:topology:interior}{partie dense} dénombrable)
possède une \omterm{def:b3:hilbert:onb}{base hilbertienne}.\index{orthonormalisation de
Gram--Schmidt}
\end{proposition}

\begin{proof}
Récurrence~: $\tilde e_n \perp e_k$ ($k < n$) par
construction, et $\tilde e_n \neq 0$ par indépendance~; les
enveloppes coïncident à chaque étape (changement de \omterm{def:b3:topology:basis}{base}
triangulaire). Pour un $H$ \omterm{def:b3:galois:separable}{séparable}~: d'une suite \omterm{def:b3:topology:interior}{dense} on
extrait une sous-famille linéairement indépendante d'enveloppe
\omterm{def:b3:topology:interior}{dense} (on jette chaque vecteur dans l'enveloppe de ses
prédécesseurs --- l'enveloppe ne change pas), on
orthonormalise~: le résultat est total.
\end{proof}

\begin{theorem}[Le système trigonométrique~; Parseval enfin]
\label{thm:b3:hilbert:fourier}
Dans $L^2(\intcc{-\pi}\pi)$ avec $\langle f, g\rangle =
\frac1{2\pi}\int_{-\pi}^\pi \bar fg$, la famille $e_n(t) =
\eu^{\iu nt}$, $n \in \Z$, est une \omterm{def:b3:hilbert:onb}{base hilbertienne}. Par
suite, pour \emph{toute} $f \in L^2$ --- en particulier toute
$f$ $2\pi$-périodique \omterm{def:b3:topology:continuity}{continue} par morceaux --- avec $c_n(f)
= \frac1{2\pi}\int_{-\pi}^{\pi}f(t)\eu^{-\iu nt}\dd t$~:
\[
f = \sum_{n\in\Z}c_n(f)\,\eu^{\iu nt} \ \ \text{dans } L^2,
\qquad
\frac1{2\pi}\int_{-\pi}^{\pi}\abs f^2 =
\sum_{n\in\Z}\abs{c_n(f)}^2 .
\]
Ceci démontre, en toute généralité, l'identité de \omterm{thm:b3:hilbert:parseval}{Parseval}
que la deuxième année admettait.\index{identité de Parseval
(générale)}
\end{theorem}

\begin{proof}
L'orthonormalité est un calcul direct (deuxième année).
Totalité~: soit $f \in L^2$ orthogonale à tous les $e_n$,
c'est-à-dire que tous les coefficients de Fourier s'annulent.
Les fonctions \omterm{def:b3:topology:continuity}{continues} $2\pi$-périodiques sont \omterm{def:b3:topology:interior}{denses} dans
$L^2(\intcc{-\pi}\pi)$~: en effet $\mathcal
C_c(\intoo{-\pi}\pi)$ est \omterm{def:b3:topology:interior}{dense}
(\cref{thm:b3:lp:density}(2)) et de telles fonctions se
prolongent périodiquement et continûment. Les polynômes
trigonométriques sont $\norm\cdot_\infty$-denses parmi les
fonctions \omterm{def:b3:topology:continuity}{continues} périodiques (Stone--Weierstrass,
\cref{cor:b3:complete:weierstrass}(c)), et $\norm\cdot_2 \leq
\norm\cdot_\infty$~: les polynômes trigonométriques sont
\omterm{def:b3:topology:interior}{denses} dans $L^2$. Or $f \perp$ tout polynôme
trigonométrique, donc $f \perp$ un sous-espace \omterm{def:b3:topology:interior}{dense}~: $f
\in (\text{dense})^\perp = \{0\}$
(\cref{thm:b3:hilbert:decomposition}). Le critère (iv) du~\cref{thm:b3:hilbert:parseval} conclut~; (ii) et (iii)
déplient l'affichage (en réindexant le $\Z$ dénombrable~; la
série bidirectionnelle converge inconditionnellement --- les
sommes partielles sur toute famille exhaustive convergent,
par l'argument de queue $\ell^2$).
\end{proof}

\begin{theorem}[Lax--Milgram]\label{thm:b3:hilbert:laxmilgram}
Soit $H$ un \omterm{def:b3:hilbert:inner}{espace de Hilbert} réel et $a \colon H\times H \to
\R$ bilinéaire, \emph{\omterm{def:b3:topology:continuity}{continue}} ($\abs{a(u,v)} \leq M\norm
u\norm v$) et \emph{coercive} ($a(u, u) \geq \alpha\norm u^2$,
$\alpha > 0$). Alors pour tout $\varphi \in H'$ il existe un
unique $u \in H$ tel que\index{théorème de Lax--Milgram}
\[
a(u, v) = \varphi(v) \qquad \text{pour tout } v \in H .
\]
\end{theorem}

\begin{proof}
Pour $u$ fixé, $v \mapsto a(u, v)$ est linéaire \omterm{def:b3:topology:continuity}{continue}~:
Riesz donne un unique $Au \in H$ avec $a(u,v) = \langle Au,
v\rangle$~; $A$ est linéaire avec $\norm{Au} \leq M\norm u$
(unicité des représentants, puis majoration). Coercivité~:
$\alpha\norm u^2 \leq a(u,u) = \langle Au, u\rangle \leq
\norm{Au}\norm u$, donc $\norm{Au} \geq \alpha\norm u$~: $A$
est injective à image fermée (une suite image de Cauchy $Au_n$
force $u_n$ de Cauchy). L'image est \omterm{def:b3:topology:interior}{dense}~: $w \perp
\operatorname{im}A$ donne $0 = \langle Aw, w\rangle \geq
\alpha\norm w^2$. Fermée et \omterm{def:b3:topology:interior}{dense}~: $A$ est bijective. Pour
$\varphi$ donné, soit $f$ son représentant (Riesz) et $u =
A^{-1}f$~: $a(u, v) = \langle f, v\rangle = \varphi(v)$,
uniquement ($a(u - u', \cdot) = 0$ et coercivité).
\end{proof}

\begin{remark}\label{rem:b3:hilbert:variational}
Lorsque $a$ est symétrique, la solution de Lax--Milgram est
l'unique minimiseur de l'\emph{énergie} $J(v) =
\frac12a(v,v) - \varphi(v)$ (\cref{exo:b3:hilbert:9})~:
existence de solutions aux problèmes variationnels d'un seul
coup. Appliqué à des espaces de fonctions adéquats (les
espaces de Sobolev d'un cours ultérieur), ceci résout les
problèmes aux limites pour les équations différentielles ---
la porte d'entrée moderne aux équations aux dérivées
partielles.
\end{remark}

\section{Exercices}

\begin{exercise}[$\star$]\label{exo:b3:hilbert:1}
(a) Démontrer les identités de polarisation (réelle~:
$4\langle x, y\rangle = \norm{x+y}^2 - \norm{x-y}^2$~;
complexe~: la version à quatre termes).
(b) Montrer que $\norm\cdot_1$ sur $L^1(\intcc01)$ et
$\norm\cdot_\infty$ sur $\mathcal C(\intcc01)$ violent
l'identité du parallélogramme~: ces normes ne proviennent
d'aucun \omterm{def:b3:hilbert:inner}{produit scalaire}.
\end{exercise}

\begin{exercise}[$\star$]\label{exo:b3:hilbert:2}
Dans $H = L^2(\intcc01)$ (réel)~:
(a) calculer la projection de $f$ sur le sous-espace des
fonctions constantes, et interpréter~;
(b) calculer la projection sur $\{g : g = 0 \text{ p.p.\ sur }
\intcc0{1/2}\}$~;
(c) calculer $d\bigl(x \mapsto x,\ \operatorname{Vect}(\mathbf
1)\bigr)$.
\end{exercise}

\begin{exercise}[$\star\star$]\label{exo:b3:hilbert:3}
(a) Montrer que pour un sous-espace $F$~: $F$ \omterm{def:b3:topology:interior}{dense} $\iff$
$F^\perp = \{0\}$, et donner un exemple dans $\ell^2$ d'un
sous-espace \omterm{def:b3:topology:interior}{dense} \emph{propre} (donc $F^\perp = 0$ sans $F
= H$~: le théorème de décomposition exige vraiment que $F$
soit fermé).
(b) Montrer que si $x_n \to x$ et $y_n \to y$ en norme, alors
$\langle x_n, y_n\rangle \to \langle x, y\rangle$, et repérer
deux endroits du chapitre où cette continuité a été utilisée.
\end{exercise}

\begin{exercise}[$\star\star$]\label{exo:b3:hilbert:4}
Appliquer Gram--Schmidt à $1, x, x^2$ dans $L^2(\intcc{-1}1)$
(\omterm{def:b3:measure:lebesgueouter}{mesure de Lebesgue})~: obtenir les trois premiers polynômes de
\emph{Legendre} normalisés, et vérifier qu'ils coïncident
avec $\sqrt{n + \frac12}\,P_n$ pour les polynômes de Rodrigues
$P_n$ du~\cref{pb:b3:hilbert:1}.
\end{exercise}

\begin{exercise}[$\star\star$]\label{exo:b3:hilbert:5}
Appliquer \omterm{thm:b3:hilbert:parseval}{Parseval} (\cref{thm:b3:hilbert:fourier}) à $f(t) =
t$ et $f(t) = t^2$ sur $\intcc{-\pi}\pi$ --- désormais
légitimement pour ces fonctions (\omterm{def:b3:topology:continuity}{continues}, mais auparavant
l'identité exigeait un soin $\mathcal C^1$ par morceaux au
saut de recollement)~: retrouver
\[
\sum_{n\geq1}\frac1{n^2} = \frac{\pi^2}6,
\qquad
\sum_{n\geq1}\frac1{n^4} = \frac{\pi^4}{90} .
\]
\end{exercise}

\begin{exercise}[$\star\star$]\label{exo:b3:hilbert:6}
(a) Trouver $a \in L^2(\intcc01)$ tel que $\int_0^{1/2}f =
\langle a, f\rangle$ pour tout $f$~; calculer $\norm\varphi$
pour cette forme.
(b) Montrer que l'évaluation $f \mapsto f(\frac12)$, définie
sur le sous-espace $\mathcal C(\intcc01) \subseteq
L^2(\intcc01)$, n'est \emph{pas} \omterm{def:b3:topology:continuity}{continue} pour
$\norm\cdot_2$~: aucun représentant de Riesz n'existe
(l'évaluation n'est pas une notion $L^2$).
\end{exercise}

\begin{exercise}[$\star\star\star$]\label{exo:b3:hilbert:7}
Soit $H$ \omterm{def:b3:galois:separable}{séparable} de \omterm{def:b3:hilbert:onb}{base hilbertienne} $(e_n)$, et $(x_k)$ une
suite bornée.
(a) Montrer qu'une sous-suite converge \emph{faiblement}~: il
existe $x$ avec $\langle y, x_{k_j}\rangle \to \langle y,
x\rangle$ pour tout $y \in H$. \emph{(Extraction diagonale
sur les coefficients $\langle e_n, x_k\rangle$~; assembler $x$
via Bessel et la bornitude uniforme des normes.)}
(b) Montrer $e_n \rightharpoonup 0$ mais $\norm{e_n} = 1$~:
les limites faibles peuvent perdre de la norme. Montrer
$\norm x \leq \liminf\norm{x_{k_j}}$ en (a).
\end{exercise}

\begin{exercise}[$\star\star$]\label{exo:b3:hilbert:8}
(Adjoints) Pour $T \in \mathcal L(H)$, montrer qu'il existe un
unique $T^* \in \mathcal L(H)$ avec $\langle Tx, y\rangle =
\langle x, T^*y\rangle$ (Riesz), et $\vertiii{T^*} =
\vertiii T$. Calculer l'adjoint du décalage $S$ sur $\ell^2$,
et prouver $\ker T^* = (\operatorname{im}T)^\perp$ --- en
déduire $\overline{\operatorname{im}T} = (\ker T^*)^\perp$.
\end{exercise}

\begin{exercise}[$\star\star$]\label{exo:b3:hilbert:9}
Soit $a$ comme dans Lax--Milgram et de plus \emph{symétrique}.
Montrer que $u$ résout $a(u, \cdot) = \varphi$ ssi $u$
minimise $J(v) = \frac12a(v, v) - \varphi(v)$, et que le
minimum est atteint en exactement un point. \emph{(Compléter
le carré~: $J(u + w) - J(u) = \frac12a(w,w) \geq
\frac\alpha2\norm w^2$.)} Application~: redériver le théorème
de projection pour les sous-espaces fermés à partir de
Lax--Milgram.
\end{exercise}

\begin{exercise}[$\star\star\star$]\label{exo:b3:hilbert:10}
(Le système de Haar) Sur $\intcc01$, soit $h_{0} = \mathbf 1$,
et pour $n = 2^j + k$ ($j \geq 0$, $0 \leq k < 2^j$)~:
\[
h_n = 2^{j/2}\Bigl(\mathbf 1_{[k2^{-j},\,(k +
\frac12)2^{-j})} - \mathbf 1_{[(k+\frac12)2^{-j},\,(k+1)2^{-j})}
\Bigr).
\]
Montrer que $(h_n)_{n\geq0}$ est orthonormée dans
$L^2(\intcc01)$, et totale. \emph{(Orthogonalité~: supports
disjoints ou emboîtés~; totalité~: les enveloppes finies
contiennent toutes les fonctions en escalier dyadiques, qui
sont \omterm{def:b3:topology:interior}{denses} --- via le~\cref{thm:b3:lp:density}(1) et
l'approximation dyadique des intervalles.)} Le système de Haar
est l'ancêtre des ondelettes.
\end{exercise}

\begin{exercise}[$\star\star$]\label{exo:b3:hilbert:11}
(\omterm{thm:b3:hilbert:decomposition}{Projections orthogonales}, caractérisées) Soit $H$ un \omterm{def:b3:hilbert:inner}{espace
de Hilbert} et $P \in \mathcal L(H)$ avec $P^2 = P$, $P \neq
0$. Montrer l'équivalence de~:
(i) $P$ est la projection orthogonale sur
$\operatorname{im}P$~;
(ii) $P = P^*$ (\cref{exo:b3:hilbert:8})~;
(iii) $\vertiii P = 1$.
\emph{(Pour (iii) $\Rightarrow$ (i)~: si un $x \in
(\ker P)^\perp$ avait $Px \neq x$, considérer $x + t(Px - x)$
--- ou directement~: pour $u \in \operatorname{im}P$ et $v
\in \ker P$, développer $\norm{P(u + tv)}^2 \leq \norm{u +
tv}^2$ pour tout $t \in \R$ et conclure $\langle u, v\rangle
= 0$.)} Exhiber une projection non orthogonale sur $\R^2$ et
calculer sa norme.
\end{exercise}

\begin{exercise}[$\star\star\star$]\label{exo:b3:hilbert:12}
(Théorème ergodique de von Neumann) Soit $U \in \mathcal L(H)$
\emph{unitaire} ($U^*U = UU^* = I$), $F = \ker(U - I)$
l'espace des points fixes, $P$ la projection orthogonale sur
$F$, et $A_n = \frac1n\sum_{k=0}^{n-1}U^k$.
(a) Montrer $\ker(U - I) = \ker(U^* - I)$ \emph{(à partir de
$\norm{Ux - x}^2 = 2\norm x^2 - 2\operatorname{Re}\langle
Ux, x\rangle$ et de l'unitarité)}, et en déduire
$\overline{\operatorname{im}(U - I)} = F^\perp$.
(b) Montrer que $A_nx \to x$ pour $x \in F$, et $A_nx \to 0$
pour $x \in \operatorname{im}(U - I)$ \emph{(télescopage)},
puis pour $x \in \overline{\operatorname{im}(U - I)}$
(borne uniforme $\vertiii{A_n} \leq 1$).
(c) Conclure~: $A_nx \to Px$ pour \emph{tout} $x \in H$ ---
les moyennes temporelles convergent vers la projection sur
les invariants.
(d) Détailler pour $H = L^2(\R/\Z)$ et $Uf = f(\cdot +
\alpha)$ avec $\alpha$ irrationnel~: identifier $F$ (utiliser
les séries de Fourier, \cref{thm:b3:hilbert:fourier}) et en
déduire que $\frac1n\sum_{k<n}f(x + k\alpha) \to \int_0^1f$
dans $L^2$~: l'équidistribution $L^2$ des rotations
irrationnelles.
\end{exercise}

\section{Problème~: polynômes orthogonaux}

\begin{problem}[{Problème du week-end --- Legendre, Hermite, et
quadrature de Gauss}]\label{pb:b3:hilbert:1}
Soit $I \subseteq \R$ un intervalle et $w > 0$ un
\emph{poids} continu sur l'intérieur de $I$ tel que $\int_I
\abs t^nw(t)\dd t < \infty$ pour tout $n$~; on travaille dans
$H = L^2(I, w\,\dd\lambda)$ avec $\langle f, g\rangle =
\int_I \bar fg\,w$. Gram--Schmidt appliqué à $1, t, t^2,
\dots$ produit les \emph{\omterm{pb:b3:hilbert:1}{polynômes orthogonaux}} $(p_n)$ pour
$w$ (normalisation monique~: $p_n = t^n + \cdots$).\index{polynômes
orthogonaux}

\medskip
\noindent\textbf{Partie I --- Théorie générale.}
\begin{enumerate}
  \item Montrer que $p_n$ est orthogonal à tout polynôme de
        \omterm{def:b3:galois:extension}{degré} $< n$, et que $(p_0, \dots, p_n)$ est une \omterm{def:b3:topology:basis}{base}
        de $\R_n[t]$.
  \item (Récurrence à trois termes) Montrer qu'il existe des
        réels $a_n, b_n$ avec
        \[
        p_{n+1}(t) = (t - a_n)\,p_n(t) - b_n\,p_{n-1}(t),
        \qquad b_n = \frac{\norm{p_n}^2}{\norm{p_{n-1}}^2} >
        0 .
        \]
        \emph{(Développer $t\,p_n$ dans la \omterm{def:b3:topology:basis}{base} $(p_k)_{k \leq
        n+1}$ et tuer les coefficients par orthogonalité, en
        utilisant $\langle tp_n, p_k\rangle = \langle p_n,
        tp_k\rangle$.)}
  \item (Racines) Montrer que $p_n$ a $n$ racines
        \emph{distinctes}, toutes intérieures à $I$.
        \emph{(Soit $t_1 < \dots < t_m$ les changements de
        signe intérieurs de $p_n$~; si $m < n$, tester $p_n$
        contre $\prod_{i\leq m}(t - t_i)$ et contredire
        l'orthogonalité.)}
\end{enumerate}

\medskip
\noindent\textbf{Partie II --- Legendre ($I = \intcc{-1}1$, $w
= 1$).} On pose $P_n(t) = \frac{1}{2^nn!}\,\frac{\dd^n}{\dd
t^n}\bigl[(t^2 - 1)^n\bigr]$ (Rodrigues).
\begin{enumerate}[resume]
  \item Montrer $\deg P_n = n$ avec coefficient dominant
        $\frac{(2n)!}{2^n(n!)^2}$, et, en intégrant par
        parties $n$ fois, que $\langle P_n, Q\rangle = 0$
        pour tout polynôme $Q$ de \omterm{def:b3:galois:extension}{degré} $< n$~: les $P_n$
        sont (à normalisation près) les \omterm{pb:b3:hilbert:1}{polynômes orthogonaux}
        pour $w = 1$.
  \item Calculer $\norm{P_n}_2^2 = \frac{2}{2n+1}$
        \emph{(intégrer par parties $n$ fois contre soi-même
        et se ramener à une intégrale Bêta/Wallis,
        \cref{exo:b3:product:8})}.
  \item Montrer que les polynômes de Legendre normalisés
        forment une \omterm{def:b3:hilbert:onb}{base hilbertienne} de $L^2(\intcc{-1}1)$
        \emph{(Weierstrass, \cref{cor:b3:complete:weierstrass},
        plus \omterm{ex:b3:lebesgue:gamma}{densité} de $\mathcal C$ dans $L^2$)}, et
        développer $f(t) = \abs t$ jusqu'au \omterm{def:b3:galois:extension}{degré} $2$~:
        calculer la meilleure approximation quadratique $L^2$
        de $\abs t$.
\end{enumerate}

\medskip
\noindent\textbf{Partie III --- Hermite ($I = \R$, $w(t) =
\eu^{-t^2}$).} On pose $H_n(t) =
(-1)^n\eu^{t^2}\frac{\dd^n}{\dd t^n}\eu^{-t^2}$.
\begin{enumerate}[resume]
  \item Montrer que $H_n$ est un polynôme de \omterm{def:b3:galois:extension}{degré} $n$ de
        coefficient dominant $2^n$, que $H_{n+1} = 2tH_n -
        H_n'$, et que $\langle H_m, H_n\rangle_w =
        \delta_{mn}\,2^nn!\sqrt\pi$ \emph{(parties encore)}.
  \item Montrer que la famille d'Hermite est totale dans
        $L^2(\R, \eu^{-t^2}\dd t)$, en admettant un résultat
        du~\cref{ch:b3:fouriertransform}~: si $g \in L^1(\R)$
        a $\int g(t)\eu^{-\iu\xi t}\dd t = 0$ pour tout $\xi$,
        alors $g = 0$ p.p. \emph{(Pour $f \perp$ tous les
        $H_n$, i.e.\ $\perp$ tous les polynômes~: montrer que
        $z \mapsto \int f(t)\eu^{-t^2}\eu^{-\iu zt}\dd t$ est
        bien définie, développer l'exponentielle en série,
        justifier l'interversion par domination, et conclure
        que la transformée de Fourier de $f\eu^{-t^2}$
        s'annule.)}
\end{enumerate}

\medskip
\noindent\textbf{Partie IV --- Quadrature de Gauss.} On fixe
$n$, on note $t_1 < \dots < t_n$ les racines de $p_n$ (Partie
I), et on définit les poids $w_i = \int_I \ell_i(t)\,w(t)\dd
t$ où les $\ell_i$ sont les polynômes de \omterm{def:b3:topology:basis}{base} de l'interpolation
de Lagrange aux $t_i$.
\begin{enumerate}[resume]
  \item Montrer que la règle de quadrature $Q(f) =
        \sum_iw_if(t_i)$ est exacte sur tous les polynômes de
        \omterm{def:b3:galois:extension}{degré} $\leq n - 1$ (interpolation), et en fait --- le
        miracle --- sur tous les polynômes de \omterm{def:b3:galois:extension}{degré} $\leq 2n
        - 1$~: écrire $P = qp_n + r$ et utiliser
        l'orthogonalité sur le quotient $q$.
  \item Montrer que les poids sont positifs \emph{(appliquer
        la règle à $\ell_i^2$, de \omterm{def:b3:galois:extension}{degré} $2n - 2$)}, et en
        déduire du théorème de Polya
        (\cref{exo:b3:banach:9}) que la quadrature de Gauss
        converge~: $Q_n(f) \to \int_I fw$ pour toute $f$
        \omterm{def:b3:topology:continuity}{continue} sur un $I$ compact.
  \item Pour $n = 2$, $I = \intcc{-1}1$, $w = 1$~: calculer
        les nœuds $\pm\frac1{\sqrt3}$ et les poids $1, 1$, et
        vérifier l'exactitude sur $1, t, t^2, t^3$ à la main.
        Comparer avec la règle des trapèzes sur les mêmes
        deux points d'évaluation.
\end{enumerate}

\medskip
\noindent\textbf{Partie V --- Chebyshev~: les polynômes qui
oscillent le mieux.} Maintenant $I = \intcc{-1}1$ et $w(t) =
\frac1{\sqrt{1 - t^2}}$.
\begin{enumerate}[resume]
  \item Montrer que $T_n(\cos\theta) = \cos n\theta$ définit
        un polynôme $T_n$ de \omterm{def:b3:galois:extension}{degré} $n$ (établir $T_{n+1} =
        2t\,T_n - T_{n-1}$ à partir d'une identité
        trigonométrique), de coefficient dominant $2^{n-1}$
        pour $n \geq 1$~; et que la substitution $t =
        \cos\theta$ donne
        \[
        \langle T_m, T_n\rangle_w =
        \int_0^\pi\cos m\theta\,\cos n\theta\,\dd\theta
        = 0 \ (m \neq n), \qquad
        \norm{T_0}_w^2 = \pi,\ \ \norm{T_n}_w^2 = \frac\pi2 :
        \]
        les $T_n$ sont les \omterm{pb:b3:hilbert:1}{polynômes orthogonaux} pour ce
        poids, et les développements de Chebyshev \emph{sont}
        des séries de Fourier en cosinus déguisées.
  \item Localiser explicitement les $n$ racines $t_k =
        \cos\frac{(2k-1)\pi}{2n}$ et les $n + 1$ extrema $s_j
        = \cos\frac{j\pi}n$ de $T_n$ sur $\intcc{-1}1$, où
        $T_n(s_j) = (-1)^j$~: le graphe \emph{équi-oscille}
        entre $\pm1$.
  \item (Minimax) Montrer que parmi tous les polynômes
        \emph{moniques} de \omterm{def:b3:galois:extension}{degré} $n$, le polynôme
        $2^{1-n}T_n$ a la plus petite norme uniforme sur
        $\intcc{-1}1$, à savoir $2^{1-n}$ --- et qu'il est
        l'unique minimiseur. \emph{(Si un monique $P$ avait
        $\sup\abs P < 2^{1-n}$, la différence $2^{1-n}T_n -
        P$, de \omterm{def:b3:galois:extension}{degré} $\leq n-1$, alternerait en signe aux $n+1$
        points d'équi-oscillation.)}
  \item Application à l'interpolation~: pour des nœuds $t_1
        < \dots < t_n$ dans $\intcc{-1}1$, l'erreur de
        l'interpolation de Lagrange d'une fonction $\mathcal
        C^n$ met en jeu $\omega(t) = \prod_i(t - t_i)$.
        Montrer que le choix des racines de Chebyshev comme
        nœuds minimise $\sup_{\intcc{-1}1}\abs\omega$, et
        donner la borne résultante $\norm{f - L_nf}_\infty
        \leq \frac{\norm{f^{(n)}}_\infty}{2^{n-1}\,n!}$ ---
        comparer avec des nœuds équidistants (énoncer le
        phénomène de Runge comme histoire d'avertissement).
  \item Vérifier $\abs{T_n'(\pm1)} = n^2$ \emph{(dériver
        $T_n(\cos\theta) = \cos n\theta$ et prendre les
        limites $\theta \to 0, \pi$)}~: des polynômes bornés
        par $1$ sur $\intcc{-1}1$ peuvent avoir une dérivée
        aussi grande que $n^2$ au bord (l'inégalité de Markov
        dit pas plus --- énoncé seulement). Où dans
        l'intervalle la borne de dérivée n'est-elle que
        $O(n)$~?
  \item (Quadrature de Chebyshev--Gauss) Montrer que la règle
        de Gauss pour le poids $w$ aux $n$ racines de
        Chebyshev a des poids \emph{égaux} $w_i = \frac\pi n$
        \emph{(exactitude sur $T_0, \dots, T_{n-1}$ plus les
        sommes trigonométriques
        $\sum_{k=1}^n\cos\bigl(j\tfrac{(2k-1)\pi}{2n}\bigr) =
        0$ pour $1 \leq j \leq n - 1$)}~: la plus uniforme de
        toutes les quadratures. L'écrire pour $n = 3$.
\end{enumerate}

\medskip
\noindent\textbf{Partie VI --- Christoffel--Darboux,
entrelacement, et matrice de Jacobi.} Retour à un poids
général~; $h_k = \norm{p_k}^2$ ($p_k$ monique), $b_k =
h_k/h_{k-1}$.
\begin{enumerate}[resume]
  \item (Norme minimale) Montrer que parmi tous les polynômes
        \emph{moniques} de \omterm{def:b3:galois:extension}{degré} $n$, l'orthogonal $p_n$ est
        l'unique de norme $L^2(w)$ minimale --- identifier
        la minimisation comme une projection orthogonale sur
        $\R_{n-1}[t]$ (\cref{thm:b3:hilbert:projection} ou
        la projection de dimension finie de deuxième année).
        La propriété minimax de la question 14 est le même
        énoncé avec $L^\infty$ à la place de $L^2$~: même
        héros, deux normes.
  \item (Christoffel--Darboux) Prouver, par récurrence sur
        $n$ via la récurrence à trois termes, l'identité
        \[
        \sum_{k=0}^{n}\frac{p_k(x)\,p_k(y)}{h_k}
        = \frac{p_{n+1}(x)\,p_n(y) -
        p_n(x)\,p_{n+1}(y)}{h_n\,(x - y)}
        \qquad (x \neq y),
        \]
        et sa forme confluente ($y \to x$)~:
        $\sum_{k\leq n}\frac{p_k(x)^2}{h_k} =
        \frac{p_{n+1}'(x)p_n(x) - p_n'(x)p_{n+1}(x)}{h_n}$.
  \item En déduire que $p_n$ et $p_{n+1}$ n'ont aucune racine
        commune, et qu'en toute racine $x_0$ de $p_{n+1}$~:
        $p_n(x_0)\,p_{n+1}'(x_0) > 0$. Conclure
        l'\emph{entrelacement} des racines~: entre deux
        racines consécutives de $p_{n+1}$ se trouve
        exactement une racine de $p_n$.
  \item (Matrice de Jacobi) Soit $J_n$ la matrice tridiagonale
        symétrique $n\times n$ de diagonale $a_0, \dots,
        a_{n-1}$ et d'entrées hors diagonale $\sqrt{b_1},
        \dots, \sqrt{b_{n-1}}$. Montrer par récurrence que
        $\det(tI_n - J_n) = p_n(t)$, donc les racines de
        $p_n$ sont les valeurs propres d'une matrice
        symétrique réelle --- redémontrant en une ligne
        qu'elles sont réelles, et (avec l'entrelacement
        ci-dessus) liant les \omterm{pb:b3:hilbert:1}{polynômes orthogonaux} au monde
        spectral du~\cref{ch:b3:spectral}.
  \item (Synthèse) Assembler le dictionnaire pour les trois
        familles classiques (Legendre, Hermite, Chebyshev)~:
        intervalle, poids, formule définissante, récurrence à
        trois termes, norme, et l'habitat naturel de chacune
        (quadrature et approximation sur les compacts~;
        analyse gaussienne~; minimax et méthodes de Fourier
        en cosinus). Une phrase sur ce que la théorie
        générale (Parties I, VI) a donné qu'aucun calcul
        individuel ne pouvait.
\end{enumerate}

\medskip
\noindent\textbf{Partie VII --- Le terme d'erreur, et le noyau
derrière les poids.} Ici $I$ est compact et $f \in \mathcal
C^{2n}(I)$.
\begin{enumerate}[resume]
  \item (Formule d'erreur de Gauss) Soit $Hf$ l'interpolant
        d'Hermite de \omterm{def:b3:galois:extension}{degré} $\leq 2n - 1$ coïncidant avec $f$
        et $f'$ aux nœuds $t_1, \dots, t_n$ (prouver son
        existence et l'erreur ponctuelle
        \[
        f(t) - Hf(t) =
        \frac{f^{(2n)}(\xi_t)}{(2n)!}\;p_n(t)^2
        \]
        par l'argument usuel de fonction auxiliaire). En
        déduire, en intégrant cette identité contre $w$ et en
        serrant entre les extrema de $f^{(2n)}$, que
        \[
        \int_I f\,w - Q_n(f)
        = \frac{f^{(2n)}(\xi)}{(2n)!}\,h_n
        \qquad\text{pour un certain } \xi \in I,
        \]
        avec $h_n = \norm{p_n}^2$ comme en Partie VI~: la
        quadrature de Gauss erre d'une $2n$-ième dérivée,
        pondérée par la norme au carré du polynôme
        orthogonal monique.
  \item (Les poids sont des valeurs de Christoffel) En
        utilisant le noyau reproduisant $K_n(x, y) =
        \sum_{k=0}^{n-1} \frac{p_k(x)p_k(y)}{h_k}$ de
        $\R_{n-1}[t]$ et l'exactitude de $Q_n$ jusqu'au
        \omterm{def:b3:galois:extension}{degré} $2n - 2$, prouver
        \[
        w_i \;=\;
        \Bigl(\,\sum_{k=0}^{n-1}
        \frac{p_k(t_i)^2}{h_k}\Bigr)^{\!-1} :
        \]
        chaque poids est la valeur en son nœud de la
        \emph{fonction de Christoffel} --- positivité des
        poids (question 10) encore, maintenant avec une
        formule exacte. Vérifier qu'elle retrouve $w_1 = w_2
        = 1$ pour $n = 2$, $I = \intcc{-1}1$, $w = 1$.
  \item (Tout se vérifie sur une intégrale) Pour le poids de
        Chebyshev et $n = 3$ nœuds, calculer les deux côtés de
        \[
        \int_{-1}^{1}\frac{t^6}{\sqrt{1 - t^2}}\,\dd t
        = \frac{5\pi}{16},
        \qquad
        Q_3(t^6) = \frac{9\pi}{32},
        \]
        donc l'erreur de quadrature est exactement
        $\frac{\pi}{32}$~; puis vérifier que la formule
        d'erreur de la question 23 prédit précisément cette
        valeur (ici $f^{(6)} = 6!$ est constante, et $h_3 =
        \norm{2^{-2}T_3}_w^2 = \frac\pi{32}$)~: théorie et
        calcul s'accordent jusqu'au dernier chiffre.
\end{enumerate}

\end{problem}
